% Unit 4 free-response set -- 2 long (10) + 3 short (4) = 32 pts.
\documentclass{apchem-exam}

\apcunit{4}
\apcunittitle{Chemical Reactions}
\apcdoctitle{Free-Response Set}
\apcsubtitle{CED 4.1--4.9 \textbullet\ 5 questions \textbullet\ 32 points \textbullet\ 50 minutes}

\begin{document}
\apctitleblock

\begin{apcnote}[Scope --- two CED exclusions apply]
Covers \ced{4.1} through \ced{4.9}, worth
\SI{7}{\percent}--\SI{9}{\percent} of the exam.

Topic \ced{4.7} excludes the terms \textbf{``oxidizing agent'' and
``reducing agent''}, and excludes \textbf{rote memorization of solubility
rules} beyond the short list in 4.7.A.5. Topic \ced{4.8} excludes
\textbf{Lewis acid--base concepts}; Br\o{}nsted--Lowry is the model the
exam uses.

You will be asked which species is oxidized and which is reduced --- never
to name an ``agent''.
\end{apcnote}

\examsection{II}{Free Response}{5 questions \textbullet\ 32 points \textbullet\ 50 minutes}{\calculatorok}

% =====================================================================
\frq{long}{10}{\ced{4.2} \ced{4.3} \ced{4.5} \ced{4.7} \textbullet\ precipitation and solution stoichiometry}

\SI{40.0}{\milli\litre} of \SI{0.150}{\Molar} \ce{Pb(NO3)2} is mixed with
\SI{50.0}{\milli\litre} of \SI{0.200}{\Molar} \ce{KI}. A bright yellow
precipitate of \ce{PbI2} forms ($M = \SI{461.0}{\gram\per\mole}$). Assume
volumes are additive.

\begin{parts}
  \item Write the balanced net ionic equation.
        \answerlines[2]{\ce{Pb^2+(aq) + 2I^-(aq) -> PbI2(s)}}
  \item Determine the limiting reactant. \workspace[3.2cm]
        \keyonly{\begin{apcanswer}$n(\ce{Pb^2+}) = 0.0400 \times 0.150 =
        \num{0.00600}$~mol.

        $n(\ce{I-}) = 0.0500 \times 0.200 = \num{0.0100}$~mol.

        The net ionic equation needs \textbf{two} iodide per lead, so
        consuming all the \ce{Pb^2+} would require
        $2(0.00600) = \num{0.0120}$~mol of iodide --- more than the
        \num{0.0100} available.

        \textbf{Iodide is limiting.}\end{apcanswer}}
  \item Calculate the mass of \ce{PbI2} formed. \workspace[2.6cm]
        \keyonly{\begin{apcanswer}Iodide is limiting and the ratio is
        $2\,\ce{I-} : 1\,\ce{PbI2}$:

        $n(\ce{PbI2}) = 0.0100/2 = \num{0.00500}$~mol

        $m = 0.00500 \times 461.0 = \mathbf{2.31}$~g\end{apcanswer}}
  \item Calculate the concentration of lead(II) ion remaining in solution.
        \answerlines[3]{Lead consumed $= \num{0.00500}$~mol, so
        $0.00600 - 0.00500 = \num{0.00100}$~mol remains.

        Total volume $= 40.0 + 50.0 = \SI{90.0}{\milli\litre}$.

        $[\ce{Pb^2+}] = 0.00100/0.0900 = \mathbf{0.0111}$~M}
  \item Identify the spectator ions and explain what makes them
        spectators. \answerlines[2]{\ce{K+} and \ce{NO3-}. Both are
        present, dissolved and unchanged, before and after the reaction,
        so they take no part in the chemical change and cancel from the
        equation.}
\end{parts}

\begin{rubric}
  \rpoint{2}{(a) 2 pts the net ionic equation with the $2$ on iodide and
    the solid state shown; 1 pt if written in full molecular form}
  \rpoint{3}{(b) 1 pt both mole quantities; 1 pt applying the $2:1$ ratio
    rather than comparing moles directly; 1 pt iodide. Comparing
    \num{0.00600} with \num{0.0100} and answering \ce{Pb^2+} earns 1}
  \rpoint{2}{(c) 1 pt dividing by 2 for the ratio; 1 pt \SI{2.31}{\gram}}
  \rpoint{2}{(d) 1 pt \num{0.00100}~mol remaining; 1 pt dividing by the
    \emph{combined} \SI{90.0}{\milli\litre}}
  \rpoint{1}{(e) both spectators named, with a reason}
\end{rubric}

% =====================================================================
\frq{long}{10}{\ced{4.9} \ced{4.6} \textbullet\ redox and a dichromate titration}

Acidified dichromate oxidizes iron(II) in solution:
\[ \ce{Cr2O7^2-(aq) + Fe^2+(aq) + H+(aq) -> Cr^3+(aq) + Fe^3+(aq) + H2O(l)}
   \quad \text{(unbalanced)} \]

\begin{parts}
  \item Assign the oxidation number of chromium on each side and state
        whether it is oxidized or reduced.
        \answerlines[2]{In \ce{Cr2O7^2-}: $2x + 7(-2) = -2$, so
        $x = \mathbf{+6}$. In \ce{Cr^3+}: $\mathbf{+3}$. The oxidation
        number falls, so chromium is \textbf{reduced}.}
  \item Balance the equation in acidic solution. \workspace[3.6cm]
        \keyonly{\begin{apcanswer}Each chromium gains 3 electrons and
        there are two of them, so the dichromate half-reaction takes
        \textbf{6} electrons. Iron loses 1 each, so six \ce{Fe^2+} are
        needed:

        \ce{Cr2O7^2- + 6Fe^2+ + 14H+ -> 2Cr^3+ + 6Fe^3+ + 7H2O}

        Charge check: left $-2 + 12 + 14 = +24$; right
        $+6 + 18 = +24$. \checkmark\end{apcanswer}}
  \item A \SI{20.00}{\milli\litre} sample of the iron(II) solution
        required \SI{24.50}{\milli\litre} of \SI{0.0150}{\Molar}
        \ce{K2Cr2O7} to reach the endpoint. Calculate $[\ce{Fe^2+}]$.
        \workspace[3.2cm]
        \keyonly{\begin{apcanswer}$n(\ce{Cr2O7^2-}) = 0.02450 \times
        0.0150 = \num{3.675e-4}$~mol.

        The balanced ratio is $1 : 6$:

        $n(\ce{Fe^2+}) = 6 \times 3.675\times10^{-4} =
        \num{2.205e-3}$~mol

        $[\ce{Fe^2+}] = 2.205\times10^{-3}/0.02000 =
        \mathbf{0.110}$~M\end{apcanswer}}
  \item State what would happen to the calculated $[\ce{Fe^2+}]$ if the
        student overshot the endpoint, and why.
        \answerlines[2]{It would come out \textbf{too high}. An overshoot
        records a larger titre than the true equivalence volume, so more
        dichromate is thought to have reacted and the calculated iron(II)
        amount is correspondingly inflated.}
\end{parts}

\begin{rubric}
  \rpoint{2}{(a) 1 pt both oxidation numbers; 1 pt ``reduced''. Per the
    \ced{4.7} exclusion, do \textbf{not} require the phrase ``oxidizing
    agent''}
  \rpoint{3}{(b) 1 pt 6 electrons on the dichromate side; 1 pt the factor
    of 6 on iron; 1 pt \ce{H+} and \ce{H2O} balanced (14 and 7). Charges
    not balancing caps this at 2}
  \rpoint{3}{(c) 1 pt moles of dichromate; 1 pt the $1:6$ ratio; 1 pt
    \SI{0.110}{\Molar}. Omitting the ratio gives \num{0.0184} and earns 1}
  \rpoint{2}{(d) 1 pt ``too high''; 1 pt because the recorded volume
    exceeds the true equivalence volume}
\end{rubric}

% =====================================================================
\frq{short}{4}{\ced{4.4} \textbullet\ physical versus chemical change}

\begin{parts}
  \item Classify each as physical or chemical, with a one-line reason:
        (i) ice melting; (ii) iron rusting; (iii) salt dissolving in
        water. \answerlines[3]{(i) \textbf{Physical} --- the molecules are
        unchanged; only the arrangement and the intermolecular forces
        between them change. (ii) \textbf{Chemical} --- new substances
        with new bonds form. (iii) \textbf{Physical} --- the ions separate
        and become hydrated, but \ce{Na+} and \ce{Cl-} themselves are
        unchanged.}
  \item Dissolving is sometimes argued to be chemical because ion--dipole
        bonds form. State the strongest argument that it is physical.
        \answerlines[2]{No new chemical substance is produced: the ions
        that were in the lattice are the same ions in solution, and
        evaporating the water recovers the original solid unchanged.}
\end{parts}

\begin{rubric}
  \rpoint{3}{1 pt each classification, each needing its reason}
  \rpoint{1}{(b) recoverability / no new substance formed}
\end{rubric}

% =====================================================================
\frq{short}{4}{\ced{4.6} \ced{4.8} \textbullet\ acid--base titration}

A \SI{25.00}{\milli\litre} sample of a monoprotic acid \ce{HA} is titrated
with \SI{0.100}{\Molar} \ce{NaOH}. The endpoint occurs at
\SI{28.75}{\milli\litre}.

\begin{parts}
  \item Calculate the concentration of the acid.
        \answerlines[2]{$n(\ce{OH-}) = 0.02875 \times 0.100 =
        \num{0.00288}$~mol, which equals $n(\ce{HA})$ at equivalence.

        $[\ce{HA}] = 0.00288/0.02500 = \mathbf{0.115}$~M}
  \item Write the net ionic equation for the reaction, given that \ce{HA}
        is a \emph{weak} acid. \answerlines[3]{\ce{HA(aq) + OH^-(aq) ->
        A^-(aq) + H2O(l)}

        A weak acid is only slightly ionized, so it must be written as the
        intact molecule \ce{HA} rather than as \ce{H+}. Only \ce{Na+} is a
        spectator.}
\end{parts}

\begin{rubric}
  \rpoint{2}{(a) 1 pt moles of base equal moles of acid at equivalence;
    1 pt \SI{0.115}{\Molar}}
  \rpoint{2}{(b) 1 pt correct equation; 1 pt \ce{HA} kept intact because
    it is weak. Writing \ce{H+ + OH^- -> H2O} earns 0 for the second point}
\end{rubric}

% =====================================================================
\frq{short}{4}{\ced{4.1} \ced{4.3} \textbullet\ equations and particulate representations}

\begin{parts}
  \item Write the balanced equation for the complete combustion of ethanol,
        \ce{C2H5OH}. \answerlines[2]{\ce{C2H5OH(l) + 3O2(g) -> 2CO2(g) +
        3H2O(g)}}
  \item A particulate diagram of the reactants shows 4 ethanol molecules
        and 9 oxygen molecules. Determine how many molecules of each
        species the product diagram should show. \workspace[2.6cm]
        \keyonly{\begin{apcanswer}The ratio needed is $1 : 3$, so 4
        ethanol would require 12 oxygen. Only 9 are present, so
        \textbf{oxygen is limiting}: it supports $9/3 = 3$ ethanol.

        Products: $3 \times 2 = \mathbf{6}$ \ce{CO2} and
        $3 \times 3 = \mathbf{9}$ \ce{H2O}, together with
        $4 - 3 = \mathbf{1}$ unreacted ethanol molecule --- which must be
        drawn, since it is still there.\end{apcanswer}}
\end{parts}

\begin{rubric}
  \rpoint{2}{(a) coefficients 1, 3, 2, 3}
  \rpoint{2}{(b) 1 pt identifying oxygen as limiting; 1 pt 6 \ce{CO2},
    9 \ce{H2O} \emph{and} the leftover ethanol. Omitting the excess
    reactant is the standard error and costs the point}
\end{rubric}

\examtotal

\end{document}
